\chapter{Osteoarthritis}
\index{Osteoarthritis}
\label{sec:Osteoarthritis}

\section{Overview}
Osteoarthritis (OA) is the most common joint disease worldwide, characterized by progressive articular cartilage degradation, subchondral bone remodeling, osteophyte formation, and synovial inflammation. It affects over 500 million people globally, with prevalence increasing with age. This autoimmune condition occurs when the immune system mistakenly attacks the body's own tissues. It follows a chronic course with periods of flare and remission. This degenerative condition involves progressive deterioration of affected tissues, leading to gradual loss of function. It is more common with advancing age. This condition affects a significant number of people worldwide and can range from mild to severe in presentation. Early recognition and appropriate management are essential for optimal outcomes. A thorough understanding of the underlying mechanisms guides treatment decisions. Patient education and self-management play important roles in long-term care.
\section{Causes}
The underlying mechanisms involve complex interactions between genetic predisposition and environmental factors. Disruption of normal physiological processes leads to the characteristic features of the condition. Multiple contributing factors may act synergistically to trigger or worsen the condition. Understanding the specific cause helps guide targeted treatment approaches.
\section{Risk Factors}
This condition happens because of mechanical and metabolic factors leading to cartilage breakdown, with risk factors including advancing age, obesity, joint injury, female sex, and genetic predisposition.

\begin{itemize}[noitemsep]
  \item Family history of autoimmune disease
  \item Female sex (higher prevalence)
  \item Specific HLA genotypes
  \item Environmental triggers (infections, smoking)
  \item Hormonal factors
  \item Certain medications
  \item Advancing age
  \item Genetic predisposition
  \item Repetitive use or overuse (joints)
  \item Previous injuries
  \item Obesity
  \item Smoking and alcohol use
  \item Genetic predisposition and family history
  \item Lifestyle factors (diet, exercise, smoking, alcohol)
  \item Environmental exposures
  \item Underlying medical conditions
  \item Medication use
\end{itemize}
\section{Signs and Symptoms}
\begin{itemize}[noitemsep]
  \item You may experience joint pain (worse with activity and relieved by rest), stiffness (morning stiffness <30 minutes, unlike RA), crepitus, reduced range of motion, and joint enlargement (osteophytes, bony enlargement). Weight-bearing joints (knees, hips) and the hands (distal interphalangeal joints---Heberden nodes, proximal interphalangeal joints---Bouchard nodes, first carpometacarpal joints) are most commonly affected. Diagnosis is primarily clinical
  \item X-rays show joint space narrowing, subchondral sclerosis, osteophytes, and subchondral cysts.
  \item Pain or discomfort in the affected area
  \item Functional impairment affecting daily activities
  \item Changes in appearance or function of affected tissues
  \item Systemic symptoms may include fatigue, fever, or weight changes
  \item Symptoms may develop gradually or appear suddenly
\end{itemize}
\section{Progression and Stages}
The condition may progress through different stages of severity over time. Early stages often involve mild or subtle symptoms that may be overlooked. Moderate stages involve more noticeable symptoms affecting daily function. Advanced stages may involve significant impairment and complications.

\begin{itemize}[noitemsep]
  \item Treatment options include patient education, weight loss, exercise, physical therapy, analgesics (acetaminophen, NSAIDs), intra-articular corticosteroid or hyaluronic acid injections, and joint replacement surgery (total knee or hip arthroplasty) for advanced disease. The outlook varies from person to person
  \item Joint replacement provides excellent outcomes for advanced disease.
\end{itemize}
\section{Complications}
\begin{itemize}[noitemsep]
  \item Disease progression leading to worsening symptoms
  \item Secondary infections or injuries
  \item Side effects of treatment
  \item Reduced quality of life
  \item Emotional and psychological impact
\end{itemize}
\section{Diagnosis}
Comprehensive medical history and physical examination Laboratory tests to assess organ function and rule out other conditions Imaging studies to visualize affected structures Specialized testing depending on the affected system Consultation with relevant specialists Monitoring of disease progression over time

\begin{itemize}[noitemsep]
  \item Comprehensive medical history and physical examination
  \item Laboratory tests to assess organ function
  \item Imaging studies to visualize affected structures
  \item Specialized testing depending on the affected system
  \item Consultation with relevant specialists
\end{itemize}
\section{Prevention}
\begin{itemize}[noitemsep]
  \item Maintain a healthy lifestyle with balanced diet and exercise
  \item Avoid known risk factors
  \item Attend regular medical check-ups for early detection
  \item Follow recommended screening guidelines
\end{itemize}
\section{Treatment Options}
Medications to manage symptoms and address underlying mechanisms Lifestyle modifications and self-care measures Physical therapy and rehabilitation Surgical intervention when indicated Regular monitoring and follow-up care Patient education and support services

\begin{itemize}[noitemsep]
  \item Medications to manage symptoms and address underlying mechanisms
  \item Lifestyle modifications and self-care measures
  \item Physical therapy and rehabilitation
  \item Surgical intervention when indicated
  \item Regular monitoring and follow-up care
\end{itemize}
\section{Living with It}
\begin{itemize}[noitemsep]
  \item Follow your treatment plan as prescribed
  \item Maintain regular follow-up with your healthcare provider
  \item Make necessary lifestyle adjustments
  \item Seek support from family, friends, or support groups
  \item Stay informed about your condition
\end{itemize}
\section{Outlook}
The outlook varies depending on the specific condition, severity at diagnosis, and response to treatment. Early intervention generally leads to better outcomes. Many conditions can be effectively managed with appropriate treatment. Regular follow-up care is important for monitoring and treatment adjustment.
